\chapter{Ethical Considerations}

% This project raises ethical, legal, and professional considerations related to the use of human feedback, domain sensitivity, and the downstream use of automated prompt optimisation systems.
% \newline

% A central ethical issue concerns human judgement and bias. Human feedback may reflect individual biases or inconsistent interpretations of quality, particularly in subjective or open-ended tasks. To mitigate these risks, the evaluation strategy could use clear evaluation criteria or blind comparisons where feasible. Any conclusions drawn from human annotations explicitly acknowledge their inherent subjectivity and potential variability.
% \newline

% There are also domain-specific risks if the framework is applied in sensitive areas such as healthcare, law, or finance. Improvements in prompt quality may increase the fluency or persuasiveness of model outputs without necessarily improving their factual accuracy or reliability. For this reason, the use of the framework in high-stakes domains should always involve appropriate domain expertise and oversight.
% \newline

% From a data protection and privacy perspective, execution traces and associated feedback may contain sensitive information in real-world deployments. While this project relies exclusively on public or synthetic datasets, any future application of the framework should ensure proper anonymisation, secure storage of traces and annotations, and compliance with relevant data protection regulations and professional standards.
